\documentclass[12pt,a4paper]{article}
% \documentclass{ctexart}

% ==================== 基础宏包 ====================
\usepackage[UTF8]{ctex}
\usepackage[margin=2.2cm]{geometry}
\usepackage{amsmath,amssymb,amsthm,mathtools}
\usepackage{array}
\usepackage{booktabs}
\usepackage{enumitem}
\usepackage{xcolor}
\usepackage{graphicx}
\usepackage{fancyhdr}
\usepackage{lastpage}
\usepackage{hyperref}
\usepackage[most]{tcolorbox}
\usepackage{titlesec}
 \usepackage{tabularx}

\hypersetup{
  hidelinks
}

% ==================== 颜色定义 ====================
\definecolor{mainblue}{RGB}{34,72,112}
\definecolor{lightblue}{RGB}{241,246,252}
\definecolor{linegray}{RGB}{210,218,228}
\definecolor{textgray}{RGB}{90,98,108}

% ==================== 页面设置 ====================
\setlength{\parindent}{2em}
\setlength{\parskip}{0.35em}
\setlength{\headheight}{25pt}
\linespread{1.2}

\pagestyle{fancy}
\fancyhf{}
\lhead{\textcolor{textgray}{\small 课程作业}}
\rhead{\textcolor{textgray}{\small \thepage/\pageref*{LastPage}}}
\cfoot{}
\renewcommand{\headrulewidth}{0.4pt}
\renewcommand{\headrule}{\hbox to\headwidth{\color{linegray}\leaders\hrule height \headrulewidth\hfill}}

% ==================== 标题样式 ====================
\titleformat{\section}
  {\Large\bfseries\color{mainblue}}
  {\thesection}
  {0.6em}
  {}

\titleformat{\subsection}
  {\large\bfseries\color{mainblue}}
  {\thesubsection}
  {0.6em}
  {}

\titlespacing*{\section}{0pt}{1.2em}{0.6em}
\titlespacing*{\subsection}{0pt}{0.8em}{0.4em}

% ==================== 自定义环境 ====================
\newtcolorbox{infobox}{
  colback=lightblue,
  colframe=mainblue,
  boxrule=0.6pt,
  arc=2mm,
  left=1.2em,
  right=1.2em,
  top=0.8em,
  bottom=0.8em
}

\newtcolorbox{answerbox}{
  colback=white,
  colframe=linegray,
  boxrule=0.5pt,
  arc=2mm,
  left=1em,
  right=1em,
  top=0.8em,
  bottom=0.8em,
  breakable
}

\newcommand{\course}{数理统计}
\newcommand{\homeworktitle}{第6次作业}
\newcommand{\studentname}{倪伟丰}
\newcommand{\studentid}{2024110089}
\newcommand{\classname}{24 大数据 2 班}
\newcommand{\teacher}{袁洪松}
\newcommand{\duedate}{2026年5月29日}

% ==================== 正文开始 ====================
\begin{document}

\begin{center}
  {\Huge\bfseries\color{mainblue}\homeworktitle}\\[0.6em]
  {\Large\course}\\[1.2em]
  {\color{linegray}\rule{0.9\textwidth}{0.8pt}}
\end{center}

\begin{infobox}
\begin{tabularx}{\textwidth}{>{\bfseries}p{2.2cm} X >{\bfseries}p{2.2cm} X}
姓名 & \studentname & 学号 & \studentid \\
班级 & \classname & 教师 & \teacher \\
提交日期 & \duedate &  &  \\
\end{tabularx}
\end{infobox}

\section{Problem 1}

The Pareto distribution is a frequently used model in the study of incomes and has the distribution function
\[
F(x;\theta_1,\theta_2)=
\begin{cases}
1-(\theta_1/x)^{\theta_2}, & x\geq \theta_1,\\
0, & \text{elsewhere},
\end{cases}
\]
where $\theta_1>0$ and $\theta_2>0$. If $X_1,X_2,\ldots,X_n$ is a random sample from this distribution, find the maximum likelihood estimators of $\theta_1$ and $\theta_2$.

\begin{answerbox}
\textbf{解答：}

由于 $X_1,\ldots,X_n$ 是来自该 Pareto 分布的随机样本，以下联合密度可写成各边际密度的乘积。由分布函数求导可得
\[
f(x;\boldsymbol\theta)=
\begin{cases}
\theta_2\theta_1^{\theta_2}x^{-(\theta_2+1)}, & x\geq \theta_1,\\
0, & \text{其他}.
\end{cases}
\]
因此似然函数为
\[
L(\theta_1,\theta_2)
=\theta_2^n\theta_1^{n\theta_2}\prod_{i=1}^n x_i^{-(\theta_2+1)}
 I_{\{\theta_1\leq x_{(1)}\}},
\]
其中 $x_{(1)}=\min\{x_1,\ldots,x_n\}$。当 $0<\theta_1\leq x_{(1)}$ 时，
\[
\ell(\theta_1,\theta_2)
=n\ln\theta_2+n\theta_2\ln\theta_1-(\theta_2+1)\sum_{i=1}^n\ln x_i.
\]
在固定 $\theta_2>0$ 时，$\ell$ 关于 $\theta_1$ 单调递增，而 $\theta_1\leq x_{(1)}$，所以
\[
\hat\theta_1=X_{(1)}.
\]
再对 $\theta_2$ 求偏导，得到
\[
\frac{\partial \ell}{\partial \theta_2}
=\frac{n}{\theta_2}+n\ln\theta_1-\sum_{i=1}^n\ln x_i.
\]
令其为 $0$，并代入 $\hat\theta_1=X_{(1)}$，得到
\[
\hat\theta_2
=\frac{n}{\sum_{i=1}^n\ln x_i-n\ln X_{(1)}}
\]
\end{answerbox}

\section{Problem 2}

假设 $X_1,\cdots,X_n$ 为取自正态分布 $N(\mu,\sigma^2)$ 的样本，其中 $\mu,\sigma$ 均未知。试计算 $\sigma$ 的无偏估计的 RCB。

\begin{answerbox}
\textbf{解答:}

正态分布的 pdf 为：

$$
f(x; \mu, \sigma) = \frac{1}{\sqrt{2\pi}\sigma} \exp\left(-\frac{(x-\mu)^2}{2\sigma^2}\right).
$$

于是：

$$
\ln f(x; \mu, \sigma) = -\frac{1}{2}\ln(2\pi) - \ln \sigma - \frac{(x-\mu)^2}{2\sigma^2}.
$$

计算 $\ln f$ 的梯度：

$$
\nabla \ln f(x; \mu, \sigma) = \left( \frac{x-\mu}{\sigma^2}, -\frac{1}{\sigma} + \frac{(x-\mu)^2}{\sigma^3} \right).
$$

于是可以计算 Fisher 信息矩阵：

$$
\mathbf{I}(\mathbf{\theta}) = \begin{pmatrix}
E\left[\left(\frac{\partial \ln f}{\partial \mu}\right)^2\right] & E\left[\frac{\partial \ln f}{\partial \mu} \frac{\partial \ln f}{\partial \sigma}\right] \\
E\left[\frac{\partial \ln f}{\partial \sigma} \frac{\partial \ln f}{\partial \mu}\right] & E\left[\left(\frac{\partial \ln f}{\partial \sigma}\right)^2\right]
\end{pmatrix}
$$

其中：

$$
\mathbf{I}(\mathbf{\theta})_{\mu\mu} = Var\left(\frac{X-\mu}{\sigma^2}\right) = \frac{1}{\sigma^2},
$$

$$
\mathbf{I}(\mathbf{\theta})_{\mu\sigma} = \mathbf{I}(\mathbf{\theta})_{\sigma\mu} = -E\left[\frac{\partial^2 \ln f}{\partial \mu \partial \sigma}\right] = -E\left[-\frac{2(X-\mu)}{\sigma^3}\right] = 0,
$$

$$
\mathbf{I}(\mathbf{\theta})_{\sigma\sigma} = E\left[\left(\frac{\partial \ln f}{\partial \sigma}\right)^2\right] = E\left[\frac{1}{\sigma^2}\left(-1+\left(\frac{X-\mu}{\sigma}\right)^2\right)^2\right]
$$

其中，$\left(\frac{X-\mu}{\sigma}\right)^2 \sim \chi^2(1)$，所以 $E\left[\left(\frac{X-\mu}{\sigma}\right)^2\right] = 1$，$Var\left(\left(\frac{X-\mu}{\sigma}\right)^2\right) = 2$。因此：

$$
\begin{aligned}
\mathbf{I}(\mathbf{\theta})_{\sigma\sigma} =& \frac{1}{\sigma^2}E\left[1-2\left(\frac{X-\mu}{\sigma}\right)^2+\left(\frac{X-\mu}{\sigma}\right)^4\right] \\
=& \frac{1}{\sigma^2}\left(1 - 2 + (2 + 1)\right) = \frac{2}{\sigma^2}. 
\end{aligned}
$$

因此，Fisher 信息矩阵为：

$$
\mathbf{I}(\mathbf{\theta}) = \begin{pmatrix}
\frac{1}{\sigma^2} & 0 \\
0 & \frac{2}{\sigma^2}
\end{pmatrix}
$$

$\sigma$ 的无偏估计量 $Y$ 的 RCB 可以计算得到：

$$
\text{Var}(Y) \geq \frac{1}{n} \mathbf{I}^{-1}(\mathbf{\theta})_{\sigma\sigma} = \frac{\sigma^2}{2n}
$$
\end{answerbox}

\section{Problem 3}

假设 $X_1,\cdots,X_n$ 为取自正态分布 $N(\mu,v)$ 的样本，其中 $\mu,v$ 均未知。考虑检验问题
\[
H_0:v=v_0 \quad \mathrm{vs}\quad H_1:v\neq v_0.
\]

\subsection{}

计算此问题的似然比 $\Lambda$。


\begin{answerbox}
\textbf{解答：}

由于 $X_1,\ldots,X_n$ 为正态样本，默认独立同分布。写出似然函数：
\[
L(X_1, X_2, \ldots, X_n; \mu, v) = \prod_{i=1}^n \frac{1}{\sqrt{2\pi v}} e^{-\frac{(X_i-\mu)^2}{2v}} = (2\pi v)^{-\frac{n}{2}} e^{-\frac{1}{2v} \sum_{i=1}^n (X_i-\mu)^2}.
\]
对数似然函数为
\[
\ell(\mu,v)=-\frac n2\ln(2\pi)-\frac n2\ln v-
\frac{1}{2v}\sum_{i=1}^n(X_i-\mu)^2.
\]

在 $\mu$ 与 $v$ 都未知的情况下，计算极大似然估计：

分别对 $\mu$ 和 $v$ 求偏导数，并令其等于零：

$$
\frac{\partial \ell}{\partial \mu} = \frac{1}{v} \sum_{i=1}^n (X_i-\mu) = 0,
$$

$$
\frac{\partial \ell}{\partial v} = -\frac{n}{2v} + \frac{1}{2v^2} \sum_{i=1}^n (X_i-\mu)^2 = 0.
$$

于是得到：

\[
\hat\mu=\bar X,
\qquad
\hat v=\frac{1}{n}\sum_{i=1}^n(X_i-\bar X)^2.
\]
记
\[
Q=\sum_{i=1}^n(X_i-\bar X)^2,
\qquad B_2=\frac{Q}{n}.
\]
在原假设 $H_0:v=v_0$ 下，求 $\mu$ 的极大似然估计 $\hat{\mu}_0$：

$$
\frac{\partial \ell_0}{\partial \mu} = \frac{1}{v_0} \sum_{i=1}^n (X_i-\mu) = 0,
$$

\[
\hat\mu_0=\bar X.
\]
于是
\[
L(\hat\mu_0,v_0)=(2\pi v_0)^{-n/2}\exp\left\{-\frac{Q}{2v_0}\right\},
\]
而
\[
L(\hat\mu,\hat v)=(2\pi B_2)^{-n/2}\exp\left\{-\frac{Q}{2B_2}\right\}
=(2\pi B_2)^{-n/2}e^{-n/2}.
\]
所以似然比为
\[
\Lambda
=\frac{L(\hat\mu_0,v_0)}{L(\hat\mu,\hat v)}
=\left(\frac{B_2}{v_0}\right)^{n/2}
e^{\frac n2-\frac{Q}{2v_0}}
=\left(\frac{B_2}{v_0}\right)^{n/2}
e^{\frac n2\left(1-\frac{B_2}{v_0}\right)}.
\]
\end{answerbox}
\subsection{}

将拒绝域 $\{\Lambda\leq c\}$ 转化为关于一个在 $H_0$ 下已知分布的检验统计量的区域（要求是精确分布而非渐进分布），并确定显著水平 $\alpha$ 下，转化后拒绝域的临界值。

\begin{answerbox}
\textbf{解答：}

在 $H_0:v=v_0$ 成立时，
\[
T=\frac{1}{v_0}\sum_{i=1}^n(X_i-\bar X)^2\sim \chi^2(n-1).
\]
又因为
\[
B_2=\frac{1}{n}\sum_{i=1}^n(X_i-\bar X)^2=\frac{v_0}{n}T,
\]
所以似然比可以写成关于 $T$ 的函数：
\[
\Lambda(T)
=\left(\frac{T}{n}\right)^{n/2}e^{\frac{n-T}{2}}.
\]
取对数得
\[
\ln\Lambda(T)=\frac n2\ln\frac{T}{n}+\frac{n-T}{2}.
\]
对 $T$ 求导：
\[
\frac{\text{d}}{\text{d}T}\ln\Lambda(T)=\frac{n}{2T}-\frac12=\frac{n-T}{2T}.
\]

令 $\frac{\text{d}}{\text{d}T}\ln\Lambda(T) > 0$，解得：

$$
\frac{n-T}{2T} > 0 \implies T < n.
$$

因此 $\Lambda(T)$ 在 $(0,n)$ 上递增，在 $(n,+\infty)$ 上递减。故拒绝域 $\{\Lambda\leq c\}$ 可转化为
\[
\{T\leq a \text{或} T\geq b\},
\qquad 0<a<n<b.
\]

根据等尾原则：

$$
P(T \leq a) = \frac{\alpha}{2} , \quad P(T \geq b) = \frac{\alpha}{2}.
$$

于是：

$$
a = \chi^2_{1-\alpha/2}(n-1), \quad b = \chi^2_{\alpha/2}(n-1)
$$

所以，拒绝域为：

$$
\left\{T \leq \chi^2_{1-\alpha/2}(n-1) \text{或} T \geq \chi^2_{\alpha/2}(n-1)\right\}
$$

\end{answerbox}


\section{Problem 4}

假设总体分布中含有未知参数 $\theta$，样本为 $X_1,\ldots,X_n$。对下列每种情况，分别找出 $\theta$ 的一个充分统计量。

\subsection{}

$X_1,\ldots,X_n$ 具有 p.m.f.
  \[
  P(X=k)=\frac{1}{\theta},\qquad k=1,2,\ldots,\theta,
  \]
  其中 $\theta\in\mathbb{N}^+$ 为未知参数。

\begin{answerbox}
\textbf{解答：}

由于 $X_1,\ldots,X_n$ 是样本，所以独立同分布。联合 p.m.f. 为
\[
f(x_1;\theta)\cdots f(x_n;\theta)
=\prod_{i=1}^n\frac1\theta I_{\{1\leq x_i\leq \theta\}}
=\frac{1}{\theta^n}I_{\{x_{(n)}\leq \theta\}}I_{\{x_{(1)}\geq 1\}},
\]
其中 $x_{(1)}=\min\{x_1,\ldots,x_n\}$，$x_{(n)}=\max\{x_1,\ldots,x_n\}$。

根据 Neyman 分解定理，令
\[
k_1(x_{(n)};\theta)=\frac{1}{\theta^n}I_{\{x_{(n)}\leq\theta\}},
\qquad
k_2(x_1,\ldots,x_n)=I_{\{x_{(1)}\geq 1\}},
\]
则 $k_2$ 不依赖于 $\theta$。所以
\[
Y=X_{(n)}=\max\{X_1,\ldots,X_n\}
\]
是 $\theta$ 的一个充分统计量。
\end{answerbox}

\subsection{}

  $X_1,\ldots,X_n$ 具有 p.m.f.
  \[
  P(X=x)=\binom{x+r-1}{r-1}\theta^r(1-\theta)^x,\qquad x=0,1,2,\ldots,
  \]
  其中 $r>0$ 为已知参数，$0<\theta<1$ 为未知参数。

\begin{answerbox}
\textbf{解答：}

由于 $X_1,\ldots,X_n$ 是样本，独立同分布。联合 p.m.f. 为
\[
\begin{aligned}
f(x_1;\theta)\cdots f(x_n;\theta)
&=\prod_{i=1}^n\binom{x_i+r-1}{r-1}\theta^r(1-\theta)^{x_i}\\
&=\left(\prod_{i=1}^n\binom{x_i+r-1}{r-1}\right)\theta^{nr}(1-\theta)^{\sum_{i=1}^n x_i}.
\end{aligned}
\]
根据 Neyman 分解定理，令
\[
k_1\left(\sum_{i=1}^n x_i;\theta\right)
=\theta^{nr}(1-\theta)^{\sum_{i=1}^n x_i},
\qquad
k_2(x_1,\ldots,x_n)=\prod_{i=1}^n\binom{x_i+r-1}{r-1},
\]
其中 $k_2$ 不依赖于 $\theta$。所以
\[
Y=\sum_{i=1}^n X_i
\]
是 $\theta$ 的一个充分统计量。
\end{answerbox}

\subsection{}

   $X_1,\ldots,X_n$ 具有 p.d.f.
  \[
  f(x;\mu,\sigma^2)=\frac{1}{\sqrt{2\pi}\sigma x}
  \exp\left\{-\frac{(\ln x-\mu)^2}{2\sigma^2}\right\},\qquad x>0,
  \]
  其中 $\mu$ 已知，$\sigma^2>0$ 为未知参数。

\begin{answerbox}
\textbf{解答：}

由于 $X_1,\ldots,X_n$ 是样本，i.i.d。记未知参数 $v=\sigma^2>0$。联合密度为
\[
\begin{aligned}
f(x_1;v)\cdots f(x_n;v)
&=\prod_{i=1}^n\frac{1}{\sqrt{2\pi}\sqrt v\,x_i}
\exp\left\{-\frac{(\ln x_i-\mu)^2}{2v}\right\}I_{\{x_i>0\}}\\
&=(2\pi)^{-n/2}v^{-n/2}
\exp\left\{-\frac{1}{2v}\sum_{i=1}^n(\ln x_i-\mu)^2\right\}
\prod_{i=1}^n x_i^{-1}I_{\{x_i>0\}}.
\end{aligned}
\]
根据 Neyman 分解定理，令
\[
k_1\left(\sum_{i=1}^n(\ln x_i-\mu)^2;v\right)
=(2\pi)^{-n/2}v^{-n/2}
\exp\left\{-\frac{1}{2v}\sum_{i=1}^n(\ln x_i-\mu)^2\right\},
\]
\[
k_2(x_1,\ldots,x_n)=\prod_{i=1}^n x_i^{-1}I_{\{x_i>0\}},
\]
其中 $k_2$ 不依赖于 $v$。所以
\[
Y=\sum_{i=1}^n(\ln X_i-\mu)^2
\]
是 $\sigma^2$ 的一个充分统计量。
\end{answerbox}

\subsection{}

  $X_1,\ldots,X_n$ 具有 p.d.f.
  \[
  f(x;\theta)=\frac{1}{2\theta}\exp\left\{-\frac{|x|}{\theta}\right\},\qquad x\in\mathbb{R},
  \]
  其中 $\theta>0$ 为未知参数。

\begin{answerbox}
\textbf{解答：}

由于 $X_1,\ldots,X_n$ 是样本，i.i.d。联合密度为
\[
f(x_1;\theta)\cdots f(x_n;\theta)
=\prod_{i=1}^n\frac{1}{2\theta}\exp\left\{-\frac{|x_i|}{\theta}\right\}
=(2\theta)^{-n}\exp\left\{-\frac{1}{\theta}\sum_{i=1}^n|x_i|\right\}.
\]
根据 Neyman 分解定理，令
\[
k_1\left(\sum_{i=1}^n|x_i|;\theta\right)
=(2\theta)^{-n}\exp\left\{-\frac{1}{\theta}\sum_{i=1}^n|x_i|\right\},
\qquad
k_2(x_1,\ldots,x_n)=1.
\]
所以
\[
Y=\sum_{i=1}^n|X_i|
\]
是 $\theta$ 的一个充分统计量。
\end{answerbox}

\subsection{}

  $X_1,\ldots,X_n$ 具有 p.d.f.
  \[
  f(x;\theta)=\frac{\theta}{x^{\theta+1}},\qquad x>1,
  \]
  其中 $\theta>0$ 为未知参数。

\begin{answerbox}
\textbf{解答：}

由于 $X_1,\ldots,X_n$ 是样本，i.i.d。联合密度为
\[
\begin{aligned}
f(x_1;\theta)\cdots f(x_n;\theta)
&=\prod_{i=1}^n\theta x_i^{-(\theta+1)}I_{\{x_i>1\}}\\
&=\theta^n\left(\prod_{i=1}^n x_i\right)^{-\theta}
\left(\prod_{i=1}^n x_i\right)^{-1}I_{\{x_{(1)}>1\}}.
\end{aligned}
\]
根据 Neyman 分解定理，令
\[
k_1\left(\prod_{i=1}^n x_i;\theta\right)
=\theta^n\left(\prod_{i=1}^n x_i\right)^{-\theta},
\qquad
k_2(x_1,\ldots,x_n)=\left(\prod_{i=1}^n x_i\right)^{-1}I_{\{x_{(1)}>1\}},
\]
其中 $k_2$ 不依赖于 $\theta$。所以
\[
Y=\prod_{i=1}^n X_i
\]
是 $\theta$ 的一个充分统计量。等价地，由于 $\log$ 为一对一变换，$\sum_{i=1}^n\log X_i$ 也是充分统计量。
\end{answerbox}

\subsection{}

  $X_1,\ldots,X_n$ 具有 p.d.f.
  \[
  f(x;\theta)=\frac{x}{\theta^2}\exp\left(-\frac{x^2}{2\theta^2}\right),\qquad x>0,
  \]
  其中 $\theta>0$ 为未知参数。

\begin{answerbox}
\textbf{解答：}

由于 $X_1,\ldots,X_n$ 是样本，i.i.d。联合密度为
\[
\begin{aligned}
f(x_1;\theta)\cdots f(x_n;\theta)
&=\prod_{i=1}^n\frac{x_i}{\theta^2}\exp\left\{-\frac{x_i^2}{2\theta^2}\right\}I_{\{x_i>0\}}\\
&=\theta^{-2n}\exp\left\{-\frac{1}{2\theta^2}\sum_{i=1}^n x_i^2\right\}
\prod_{i=1}^n x_i I_{\{x_i>0\}}.
\end{aligned}
\]
根据 Neyman 分解定理，令
\[
k_1\left(\sum_{i=1}^n x_i^2;\theta\right)
=\theta^{-2n}\exp\left\{-\frac{1}{2\theta^2}\sum_{i=1}^n x_i^2\right\},
\]
\[
k_2(x_1,\ldots,x_n)=\prod_{i=1}^n x_i I_{\{x_i>0\}},
\]
其中 $k_2$ 不依赖于 $\theta$。所以
\[
Y=\sum_{i=1}^nX_i^2
\]
是 $\theta$ 的一个充分统计量。
\end{answerbox}

\section{Problem 5}

假设 $X_1,\cdots,X_n$ 为来自某分布的样本，其 p.d.f. 为
\[
f(x;a,b)=
\begin{cases}
\dfrac{1}{b}e^{-(x-a)/b}, & x>a,\\
0, & x\leq a.
\end{cases}
\]

\subsection{}

如果 $a=a_0$ 已知，请找出 $b$ 的一个充分统计量。

\begin{answerbox}
\textbf{解答：}

由于 $X_1,\ldots,X_n$ 是样本，以下按独立同分布处理。当 $a=a_0$ 已知时，联合密度为
\[
\begin{aligned}
f(x_1;b)\cdots f(x_n;b)
&=\prod_{i=1}^n\frac1b\exp\left\{-\frac{x_i-a_0}{b}\right\}I_{\{x_i>a_0\}}\\
&=b^{-n}\exp\left\{-\frac{1}{b}\sum_{i=1}^n(x_i-a_0)\right\}I_{\{x_{(1)}>a_0\}}.
\end{aligned}
\]
根据 Neyman 分解定理，令
\[
k_1\left(\sum_{i=1}^n(x_i-a_0);b\right)
=b^{-n}\exp\left\{-\frac{1}{b}\sum_{i=1}^n(x_i-a_0)\right\},
\]
\[
k_2(x_1,\ldots,x_n)=I_{\{x_{(1)}>a_0\}},
\]
其中 $k_2$ 不依赖于 $b$。所以
\[
Y=\sum_{i=1}^nX_i
\]
是 $b$ 的一个充分统计量。
\end{answerbox}

\subsection{}

如果 $b=b_0$ 已知，请找出 $a$ 的一个充分统计量。

\begin{answerbox}
\textbf{解答：}

由于 $X_1,\ldots,X_n$ 是样本，以下按独立同分布处理。当 $b=b_0$ 已知时，联合密度为
\[
\begin{aligned}
f(x_1;a)\cdots f(x_n;a)
&=\prod_{i=1}^n\frac{1}{b_0}\exp\left\{-\frac{x_i-a}{b_0}\right\}I_{\{x_i>a\}}\\
&=b_0^{-n}\exp\left\{-\frac{1}{b_0}\sum_{i=1}^nx_i\right\}
\exp\left\{\frac{na}{b_0}\right\}I_{\{a<x_{(1)}\}}.
\end{aligned}
\]
这里含参数 $a$ 的支持条件只能通过 $x_{(1)}$ 进入。根据 Neyman 分解定理，令
\[
k_1(x_{(1)};a)=e^{\frac{na}{b_0}}I_{\{a<x_{(1)}\}},
\]
\[
k_2(x_1,\ldots,x_n)=b_0^{-n}e^{-\frac{1}{b_0}\sum_{i=1}^nx_i},
\]
其中 $k_2$ 不依赖于 $a$。所以
\[
Y=X_{(1)}=\min\{X_1,\ldots,X_n\}
\]
是 $a$ 的一个充分统计量。
\end{answerbox}

\section{Problem 6}

假设 $X_1,\cdots,X_n$ 为来自指数分布 $\operatorname{Exp}(1/\theta)$ 的样本，其 p.d.f. 为
\[
f(x;\theta)=\frac{1}{\theta}e^{-x/\theta},\qquad x>0.
\]
我们知道，$Y_1=\sum_{i=1}^{n}X_i$ 为一个充分统计量，而 $Y_2=X_1$ 为一个无偏估计量。请用 Rao--Blackwell 定理，计算无偏估计量 $E(Y_2\mid Y_1)$。（提示：仿照课件中 Poisson 分布的做法）

\begin{answerbox}
\textbf{解答：}

由于 $X_1,\ldots,X_n$ 是来自 $\operatorname{Exp}(1/\theta)$ 的样本，故它们独立同分布，且具有轮换对称性。已知
\[
Y_1=\sum_{i=1}^nX_i
\]
为 $\theta$ 的充分统计量，且 $Y_2=X_1$ 是 $\theta$ 的无偏估计量。根据 Rao--Blackwell 定理，
\[
\varphi(Y_1)=E(Y_2\mid Y_1)
\]
仍是 $\theta$ 的无偏估计量，且方差不超过 $Y_2$ 的方差。

现在计算
\[
\varphi(y_1)=E\left(X_1\mid \sum_{i=1}^nX_i=y_1\right).
\]
由交换性，
\[
E\left(X_1\mid Y_1=y_1\right)=\cdots=E\left(X_n\mid Y_1=y_1\right).
\]
因此
\[
\sum_{i=1}^nE(X_i\mid Y_1=y_1)
=E\left(\sum_{i=1}^nX_i\mid Y_1=y_1\right)
=y_1.
\]
所以
\[
\varphi(y_1)=\frac{y_1}{n},
\qquad
\varphi(Y_1)=\frac{Y_1}{n}=\bar X.
\]
并且
\[
E(\bar X)=\theta,
\qquad
\operatorname{Var}(\bar X)=\frac{1}{n^2}\sum_{i=1}^n\operatorname{Var}(X_i)=\frac{\theta^2}{n}\leq \theta^2=\operatorname{Var}(X_1).
\]
故 Rao--Blackwell 化后的无偏估计量为 $E(Y_2\mid Y_1)=\frac{Y_1}{n}=\bar X$.
\end{answerbox}

\section{Problem 7}

设 $X_1,\cdots,X_n$ 为来自几何分布 $\operatorname{Geometric}(\theta)$ 的样本，其 p.m.f. 为
\[
P(X=x)=\theta(1-\theta)^{x-1},\qquad x=1,2,\ldots,\quad 0<\theta<1.
\]

\subsection{}

证明：$T=\sum_{i=1}^{n}X_i$ 是 $\theta$ 的充分统计量，且服从帕斯卡（负二项）分布：
  \[
  P(T=t)=\binom{t-1}{n-1}\theta^n(1-\theta)^{t-n},\qquad t=n,n+1,n+2,\ldots.
  \]
  （提示：几何分布可视为首次成功所需的试验次数，而负二项分布可视为 $n$ 次成功所需的总试验次数）

\begin{answerbox}
\textbf{解答：}

\textbf{证明 $T$ 是 $\theta$ 的充分统计量。}
由于 $X_1,\ldots,X_n$ 是样本，以下按独立同分布处理。联合 p.m.f. 为
\[
f(x_1;\theta)\cdots f(x_n;\theta)
=\prod_{i=1}^n\theta(1-\theta)^{x_i-1}
=\theta^n(1-\theta)^{\sum_{i=1}^nx_i-n}.
\]
根据 Neyman 分解定理，令
\[
k_1\left(\sum_{i=1}^nx_i;\theta\right)
=\theta^n(1-\theta)^{\sum_{i=1}^nx_i-n},
\qquad
k_2(x_1,\ldots,x_n)=1,
\]
则 $T=\sum_{i=1}^nX_i$ 是 $\theta$ 的充分统计量。

\textbf{证明 $T$ 服从帕斯卡（负二项）分布}

几何分布可视为首次成功所需的试验次数，而负二项分布可视为 $n$ 次成功所需的总试验次数。所以 $P(T=t)$ 可以理解为 在 $t-1$ 次试验中有 $n-1$ 次成功，第 $t$ 次试验成功的概率。

所以，在前 $t-1$ 次试验中成功的次数服从 $\operatorname{Binomial}(t-1, \theta)$ 分布，而第 $t$ 次试验成功的概率为 $\theta$，所以：

$$
\begin{aligned}
P(T = t) =& P(\text{在前 } t-1 \text{ 次试验中成功 } n-1 \text{ 次}) \cdot P(\text{第 } t \text{ 次试验成功}) \\
=& \binom{t-1}{n-1} \theta^{n-1} (1-\theta)^{(t-1)-(n-1)} \cdot \theta \\
=& \binom{t-1}{n-1} \theta^n (1-\theta)^{t-n}.
\end{aligned}
$$

所以，$T$ 服从帕斯卡（负二项）分布得证。

\end{answerbox}
  
\subsection{}

  计算 $E(T)$，并由此求 $\theta^{-1}$ 的一个无偏估计量。

\begin{answerbox}
\textbf{解答：}

由于 $X_i\sim\operatorname{Geometric}(\theta)$，且这里的几何分布表示首次成功所需试验次数，所以
\[
E(X_i)=\frac1\theta.
\]
又因为样本独立同分布，
\[
E(T)=E\left(\sum_{i=1}^nX_i\right)=\sum_{i=1}^nE(X_i)=\frac{n}{\theta}.
\]
因此
\[
E\left(\frac{T}{n}\right)=\frac1\theta.
\]
所以 $\frac{T}{n}=\bar X$ 是 $\theta^{-1}$ 的一个无偏估计量。
\end{answerbox}

\subsection{}

  定义统计量
  \[
  \psi(X_1)=
  \begin{cases}
  1, & X_1=1,\\
  0, & X_1\geq 2.
  \end{cases}
  \]
  试证明 $\psi(X_1)$ 是 $\theta$ 的无偏估计。当 $n\geq 2$ 时，利用 Rao--Blackwell 定理，得到 $\theta$ 的一个无偏估计量，且其方差不超过 $\psi(X_1)$ 的方差。

\begin{answerbox}
\textbf{解答：}

首先，
\[
E\{\psi(X_1)\}=1 \cdot P(X_1=1) + 0 \cdot P(X_1\geq 2) =\theta,
\]
所以 $\psi(X_1)$ 是 $\theta$ 的无偏估计量。

当 $n\geq 2$ 时，由第 (1) 问可知 $T=\sum_{i=1}^nX_i$ 是 $\theta$ 的充分统计量。又 $\psi(X_1)$ 是 $\theta$ 的无偏估计量，因此 Rao--Blackwell 定理给出
\[
\varphi(T)=E\{\psi(X_1)\mid T\}
\]
也是 $\theta$ 的无偏估计量，且
\[
\operatorname{Var}\{\varphi(T)\}\leq \operatorname{Var}\{\psi(X_1)\}.
\]
下面计算 $\varphi(T)$。对 $t=n,n+1,\ldots$，
\[
\begin{aligned}
\varphi(t)
&=E\{\psi(X_1)\mid T=t\}\\
&=P(X_1=1\mid T=t)\\
&=\frac{P(X_1=1,\,X_2+\cdots+X_n=t-1)}{P(T=t)}.
\end{aligned}
\]
由于 $X_1$ 与 $X_2,\ldots,X_n$ 独立，且 $X_2+\cdots+X_n$ 服从参数为 $(n-1,\theta)$ 的帕斯卡分布，
\[
\begin{aligned}
\varphi(t)
&=\frac{\theta\binom{t-2}{n-2}\theta^{n-1}(1-\theta)^{t-n}}
{\binom{t-1}{n-1}\theta^n(1-\theta)^{t-n}}\\
&=\frac{\binom{t-2}{n-2}}{\binom{t-1}{n-1}}
=\frac{n-1}{t-1}.
\end{aligned}
\]
因此 Rao--Blackwell 化后的估计量为
\[
\varphi(T)=\frac{n-1}{T-1}
\]
它是 $\theta$ 的无偏估计量，且方差不超过 $\psi(X_1)$ 的方差。
\end{answerbox}

\section{Problem 8}

假设 $X_1,\ldots,X_n$ 为取自 $\operatorname{Bernoulli}(\theta)$ 分布的样本。我们希望估计
\[
g(\theta)=(1-\theta)^2.
\]

\subsection{}

试找出 $g(\theta)$ 的极大似然估计量，并判断其是否无偏。

\begin{answerbox}
\textbf{解答：}

由于 $X_1,\ldots,X_n$ 为 Bernoulli 样本，i.i.d。似然函数为
\[
L(\theta)=\prod_{i=1}^n\theta^{X_i}(1-\theta)^{1-X_i}
=\theta^{\sum_{i=1}^nX_i}(1-\theta)^{n-\sum_{i=1}^nX_i}.
\]
对数似然函数为
\[
\ell(\theta)=\left(\sum_{i=1}^nX_i\right)\ln\theta+
\left(n-\sum_{i=1}^nX_i\right)\ln(1-\theta).
\]
令 $\ell'(\theta)=0$，得到
\[
\hat\theta=\bar X.
\]
由 MLE 的不变性，
\[
\hat{g}(\theta)=(1-\hat\theta)^2=(1-\bar X)^2.
\]

下面判断其是否无偏。因为 $\bar X$ 的期望与方差为
\[
E(\bar X)=\theta,
\qquad
\operatorname{Var}(\bar X)=\frac{\theta(1-\theta)}{n},
\]
所以
\[
E(\bar X^2)=\operatorname{Var}(\bar X)+\{E(\bar X)\}^2
=\frac{\theta(1-\theta)}{n}+\theta^2.
\]
因此
\[
\begin{aligned}
E\{(1-\bar X)^2\}
&=1-2E(\bar X)+E(\bar X^2)\\
&=(1-\theta)^2+\frac{\theta(1-\theta)}{n}.
\end{aligned}
\]
不等于 $(1-\theta)^2$，所以 $(1-\bar X)^2$ 不是 $g(\theta)$ 的无偏估计量。
\end{answerbox}

\subsection{}

 证明 $Y_1=\sum_{i=1}^{n}X_i$ 为 $\theta$ 的充分统计量。

\begin{answerbox}
\textbf{解答：}

由于 $X_1,\ldots,X_n$ 为 Bernoulli 样本，独立同分布。联合 p.m.f. 为
\[
f(x_1;\theta)\cdots f(x_n;\theta)
=\prod_{i=1}^n\theta^{x_i}(1-\theta)^{1-x_i}
=\theta^{\sum_{i=1}^nx_i}(1-\theta)^{n-\sum_{i=1}^nx_i}.
\]
根据 Neyman 分解定理，令
\[
k_1\left(\sum_{i=1}^nx_i;\theta\right)
=\theta^{\sum_{i=1}^nx_i}(1-\theta)^{n-\sum_{i=1}^nx_i},
\qquad
k_2(x_1,\ldots,x_n)=1.
\]
其中 $k_2$ 不依赖于 $\theta$。所以
\[
Y_1=\sum_{i=1}^nX_i
\]
是 $\theta$ 的充分统计量。
\end{answerbox}

\subsection{}

 我们考虑如下统计量：$Y_2=I_{\{X_1+X_2=0\}}$。试证明 $Y_2$ 为 $g(\theta)$ 的无偏估计。
  
\begin{answerbox}
\textbf{解答：}

当 $n\geq 2$ 时，由于 $X_1$ 与 $X_2$ 独立，
\[
E(Y_2)=P(X_1+X_2=0)=P(X_1=0,X_2=0)=P(X_1=0)P(X_2=0)=(1-\theta)^2.
\]
所以 $Y_2=I_{\{X_1+X_2=0\}}$ 是 $g(\theta)=(1-\theta)^2$ 的一个无偏估计量。
\end{answerbox}

\subsection{}

 利用 Rao--Blackwell 定理，试计算出一个 $g(\theta)$ 的无偏估计量，使它为 $Y_1$ 的函数，且有效性不低于 $Y_2$。

\begin{answerbox}
\textbf{解答：}

根据 Rao--Blackwell 定理，我们可以计算出一个 $g(\theta)$ 的无偏估计量 $\varphi(Y_1) = E(Y_2 | Y_1)$，其中 $Y_1 = \sum_{i=1}^n X_i$。

现在有：

$$
\begin{aligned}
\varphi(y_1) =& E(Y_2 | Y_1 = y_1) \\
=& P(X_1 + X_2 = 0 | Y_1 = y_1) \\
=& \frac{P\left(X_1 = 0, X_2 = 0, \sum_{i=3}^n X_i = y_1\right)}{P\left(\sum_{i=1}^n X_i = y_1\right)} \\
=& \frac{P(X_1 = 0) P(X_2 = 0) P\left(\sum_{i=3}^n X_i = y_1\right)}{P\left(\sum_{i=1}^n X_i = y_1\right)} \\
=& \frac{(1-\theta)^2 \binom{n-2}{y_1} \theta^{y_1} (1-\theta)^{(n-2)-y_1}}{\binom{n}{y_1} \theta^{y_1} (1-\theta)^{n-y_1}} \\
=& \frac{(1-\theta)^2 \binom{n-2}{y_1}}{\binom{n}{y_1} (1-\theta)^2} \\
=& \frac{(n-y_1)(n-y_1-1)}{n(n-1)}.
\end{aligned}
$$

于是，根据 Rao--Blackwell 定理，$\varphi(Y_1) = \frac{(n-Y_1)(n-Y_1-1)}{n(n-1)}$ 是 $g(\theta)$ 的一个无偏估计量，且有效性不低于 $Y_2$。


\end{answerbox}

\end{document}
